% Part: computability
% Chapter: computability-theory
% Section: applications-fixed-point

\documentclass[../../../include/open-logic-section]{subfiles}

\begin{document}

\olfileid{cmp}{thy}{apf}
\olsection{تطبيق مبرهنة النقطة الثابتة}

فحاصل مبرهنة النقطة الثابتة أنها تتيح تعريف الدوال الجزئية القابلة
للحساب بالرجوع إلى فهارسها. ومن ذلك أن نجد فهرسًا $\OLId{latin-ordinary}{e}$ يحقق، لكل $\OLId{latin-ordinary}{y}$،
\[
\cfind{\OLId{latin-ordinary}{e}}(\OLId{latin-ordinary}{y}) = \OLId{latin-ordinary}{e} + \OLId{latin-ordinary}{y}.
\]
ومن أمثلتها أيضًا أن نستعمل برهانها في إنشاء برنامج
بلغة Java أو C++ يطبع نفسه.

وقد جعلنا، لكل $\OLId{latin-ordinary}{e}$، المجموعة $\OLId{latin-looped}{W}_\OLId{latin-ordinary}{e}$ مجال $\cfind{\OLId{latin-ordinary}{e}}$؛
فكانت $\OLId{latin-looped}{W}_\OLMathNumeral{0}$، و$\OLId{latin-looped}{W}_\OLMathNumeral{1}$، و$\OLId{latin-looped}{W}_\OLMathNumeral{2}$،~\dots تعدادًا للمجموعات القابلة للتعداد
حسابيًا. ومنها مجموعات قابلة للحساب. فهل من خوارزمية
تتلقى $\OLId{latin-ordinary}{x}$، ومتى كانت $\OLId{latin-looped}{W}_\OLId{latin-ordinary}{x}$ قابلة للحساب أعطت فهرسًا
لدالتها المميزة؟ يمتنع ذلك، كما يأتي.

\begin{thm}
لا توجد دالة جزئية قابلة للحساب $\OLId{latin-ordinary}{f}$ لها الخاصية الآتية: متى كانت
$\OLId{latin-looped}{W}_\OLId{latin-ordinary}{e}$ قابلة للحساب كانت $\OLId{latin-ordinary}{f}(\OLId{latin-ordinary}{e})$ معرّفة وكانت $\cfind{\OLId{latin-ordinary}{f}(\OLId{latin-ordinary}{e})}$ دالتها
المميزة.
\end{thm}

\begin{proof}
نأخذ $\OLId{latin-ordinary}{f}$ جزئية قابلة للحساب كيفما اتفق، ونبني فهرسًا $\OLId{latin-ordinary}{e}$ تكون معه
$\OLId{latin-looped}{W}_\OLId{latin-ordinary}{e}$ قابلة للحساب، ولا تكون $\cfind{\OLId{latin-ordinary}{f}(\OLId{latin-ordinary}{e})}$ دالتها المميزة. وتتيح
مبرهنة النقطة الثابتة اختيار $\OLId{latin-ordinary}{e}$ بحيث
\[
\cfind{\OLId{latin-ordinary}{e}}(\OLId{latin-ordinary}{y}) \simeq
\begin{cases}
\OLMathNumeral{0} & \text{إذا كان $y=0$ وكانت $\cfind{f(e)}(0) \fdefined = 0$} \\
\text{غير معرّفة} & \text{خلاف ذلك.}
\end{cases}
\]
فالفهرس $\OLId{latin-ordinary}{e}$ حاصل من تطبيق مبرهنة النقطة الثابتة على الدالة ذات
التعريف
\[
\OLId{latin-ordinary}{g}(\OLId{latin-ordinary}{x},\OLId{latin-ordinary}{y}) \simeq
\begin{cases}
\OLMathNumeral{0} & \text{إذا كان $y=0$ وكانت $\cfind{f(x)}(0) \fdefined = 0$} \\
\text{غير معرّفة} & \text{خلاف ذلك.}
\end{cases}
\]
أما قابلية $\OLId{latin-ordinary}{g}$ للحساب الجزئي فتظهر من وصف الإجراء الآتي:
إذا أُدخل $\OLId{latin-ordinary}{x}$ و$\OLId{latin-ordinary}{y}$ فحصت الخوارزمية أولًا مساواة $\OLId{latin-ordinary}{y}$ لـ
$\OLMathNumeral{0}$؛ فإن تساويا حسبت $\OLId{latin-ordinary}{f}(\OLId{latin-ordinary}{x})$، ثم أجرت الآلة الشاملة
لحساب $\cfind{\OLId{latin-ordinary}{f}(\OLId{latin-ordinary}{x})}(\OLMathNumeral{0})$. فإن توقف الحساب وأعطى~$\OLMathNumeral{0}$،
أعطت الخوارزمية~$\OLMathNumeral{0}$؛ وإلا لم تتوقف.

فنقسم النظر: إن كانت $\OLId{latin-ordinary}{f}(\OLId{latin-ordinary}{e})$ غير معرّفة، كان $\OLId{latin-looped}{W}_\OLId{latin-ordinary}{e}=\emptyset$،
وفات الشرط الذي يوجب تعريف $\OLId{latin-ordinary}{f}(\OLId{latin-ordinary}{e})$. وإن كانت $\OLId{latin-ordinary}{f}(\OLId{latin-ordinary}{e})$ معرّفة، وكانت
$\cfind{\OLId{latin-ordinary}{f}(\OLId{latin-ordinary}{e})}(\OLMathNumeral{0})$ معرّفة وقيمتها $\OLMathNumeral{0}$، تعينت $\cfind{\OLId{latin-ordinary}{e}}(\OLId{latin-ordinary}{y})$
إذا وفقط إذا ساوى $\OLId{latin-ordinary}{y}$ العدد $\OLMathNumeral{0}$؛ فكان $\OLId{latin-looped}{W}_\OLId{latin-ordinary}{e}=\{\OLMathNumeral{0}\}$. وإن كانت
$\OLId{latin-ordinary}{f}(\OLId{latin-ordinary}{e})$ معرّفة، ولم تُعرّف $\cfind{\OLId{latin-ordinary}{f}(\OLId{latin-ordinary}{e})}(\OLMathNumeral{0})$ أو لم تساوِ $\OLMathNumeral{0}$، كان
$\OLId{latin-looped}{W}_\OLId{latin-ordinary}{e}=\emptyset$. فلا تكون $\cfind{\OLId{latin-ordinary}{f}(\OLId{latin-ordinary}{e})}$ في أي من الحالتين الأخيرتين دالة
مميزة لـ~$\OLId{latin-looped}{W}_\OLId{latin-ordinary}{e}$، إذ تخفق في إعطاء الجواب الصحيح عند الدخل $\OLMathNumeral{0}$.
\end{proof}

\end{document}
